\subsection{Integration with respect to a vectorial measure}

Let \(\mathbf{m}\) be a vectorial measure on \(T\) , with values in \(F\) . For every \(\mathbf{z}' \in F'\) , the mapping \(\mathbf{z}' \circ \mathbf{m}\) is a scalar measure on \(T\) , depending linearly on \(\mathbf{z}'\) . If \(f\) is a numerical function defined on \(T\) , we shall say, by an abuse of language, that the pair \((f, \mathbf{m})\) has the property \(\boldsymbol{P}\) if, for every \(\mathbf{z}' \in F'\) , the pair \((f, |\mathbf{z}' \circ \mathbf{m}|)\) has the property \(\boldsymbol{P}\) . For example, we shall say that \(f\) is essentially integrable for \(\mathbf{m}\) if, for every \(\mathbf{z}' \in F'\) , the function \(f\) is essentially integrable for \(|\mathbf{z}' \circ \mathbf{m}|\) . It comes to the same to say that \(f\) is essentially integrable for each of the measures \((\mathbf{z}' \circ \mathbf{m})^+\) and \((\mathbf{z}' \circ \mathbf{m})^-\) (Ch. V, §2, No.~2, Cor. 2 of Prop. 3).

PROPOSITION 1. --- Let m be a vectorial measure on T with values in F, f a numerical function on T that is essentially integrable for m. The mapping

\[
\mathbf {z} ^ {\prime} \mapsto \int f d (\mathbf {z} ^ {\prime} \circ \mathbf {m}) ^ {+} - \int f d (\mathbf {z} ^ {\prime} \circ \mathbf {m}) ^ {-}
\]

is a linear form on \(\mathbf{F}'\) .

Denoting this mapping by \(\Phi\) , it is immediate that \(\Phi(\lambda\mathbf{z}') = \lambda\Phi(\mathbf{z}')\) for all \(\lambda \in \mathbf{R}\) . Everything comes down to showing that \(\Phi(\mathbf{y}' + \mathbf{z}') = \Phi(\mathbf{y}') + \Phi(\mathbf{z}')\) . Set \(\mu = |\mathbf{y}' \circ \mathbf{m}| + |\mathbf{z}' \circ \mathbf{m}|\) ; by the Lebesgue-Nikodym theorem, one can then write \(\mathbf{y}' \circ \mathbf{m} = g \cdot \mu\) and \(\mathbf{z}' \circ \mathbf{m} = h \cdot \mu\) , where \(g\) and \(h\) are two bounded and locally \(\mu\) -integrable numerical functions (Ch. V, §5, No.~5, Th. 2); moreover, \((\mathbf{y}' \circ \mathbf{m})^+ = g^+ \cdot \mu\) and \((\mathbf{y}' \circ \mathbf{m})^- = g^- \cdot \mu\) , and the analogous relations hold with \(\mathbf{y}'\) replaced by \(\mathbf{z}'\) (resp. \(\mathbf{y}' + \mathbf{z}'\) ) and \(g\) by \(h\) (resp. \(g + h\) ). This being so, it is immediate that \(f\) is essentially \(\mu\) -integrable (Ch. V, §2, No.~2, Cor. 1 of Prop. 3), and the relation to be proved reduces to \((g + h)^+ - (g + h)^- = (g^+ - g^-) + (h^+ - h^-)\) , which is obvious.

DEFINITION 2. --- Let m be a vectorial measure on T with values in F, f a numerical function on T that is essentially integrable for m. One calls integral of f with respect to m, and denotes by \(\mathbf{m}(f)\) or \(\int f dm\) or again \(\int f(t) dm(t)\) , the element of \(F'^{*}\) defined by

\[
\left\langle \mathbf {z} ^ {\prime}, \int f d \mathbf {m} \right\rangle = \int f d (\mathbf {z} ^ {\prime} \circ \mathbf {m}) ^ {+} - \int f d (\mathbf {z} ^ {\prime} \circ \mathbf {m}) ^ {-}.\tag{1}
\]

We observe that if \(f \in \mathcal{K}(T)\) , the element \(\int f dm\) so defined coincides with the element denoted likewise in No.~1, because the second member of (1) is then \(\int f d(\mathbf{z}' \circ \mathbf{m}) = \mathbf{z}'(\mathbf{m}(f))\) by definition. Moreover, if in particular one applies Def. 2 to the case that F = R, one sees that for every \(z' \in F'\) , \(f\) is essentially integrable for the real measure \(\mathbf{z}' \circ \mathbf{m}\) , and that the second member of (1) may be written \(\int f d(\mathbf{z}' \circ \mathbf{m})\) .

Suppose now that \(f\) is essentially integrable for \(\mathbf{m}\) , and let \(\mathbf{z}' \in \mathbf{F}'\) . Set \(\mu = |\mathbf{z}' \circ \mathbf{m}|\) ; by the Lebesgue-Nikodym theorem, one can write \(\mathbf{z}' \circ \mathbf{m} = g \cdot \mu\) , where \(g\) is locally \(\mu\) -integrable and \(\| g \| \leqslant 1\) , and the proof of Prop. 1 shows that \(\int f d(\mathbf{z}' \circ \mathbf{m}) = \int fg d\mu\) . Consequently,

\[
\left| \int f d (\mathbf {z} ^ {\prime} \circ \mathbf {m}) \right| \leqslant \int | f | d | \mathbf {z} ^ {\prime} \circ \mathbf {m} |.\tag{2}
\]

It is clear that the set of finite numerical functions essentially integrable for m is a vector space over R; we shall denote by \(\mathcal{L}(\mathbf{m})\) this space equipped with the coarsest locally convex topology making continuous all of the linear forms \(f \mapsto \int f d(\mathbf{z}' \circ \mathbf{m})\) , where \(z'\) runs over \(F'\) . Note that in general the locally convex space \(\mathcal{L}(\mathbf{m})\) is not Hausdorff.

Example. --- Let us take for m the identity mapping of \(\mathcal{H}(\mathrm{T})\) onto itself. Since the dual of \(\mathcal{H}(\mathrm{T})\) is the space \(\mathcal{M}(\mathrm{T})\) of scalar measures on T, the functions \(f \in \mathcal{L}(\mathbf{m})\) are those that are essentially integrable for every scalar measure \(\mu\) (cf.~Exer. 1), and the integral \(\int f dm\) is the linear form \(\mu \mapsto \int f d\mu\) on \(\mathcal{M}(\mathrm{T})\) . One cannot have \(\int f d\mu = 0\) for every measure \(\mu \in \mathcal{M}(\mathrm{T})\) unless \(f = 0\) , as one sees on taking \(\mu = \varepsilon_t\) , where \(t\) is arbitrary in T; in other words, the mapping \(f \mapsto \int f dm\) is an injection of \(\mathcal{L}(\mathbf{m})\) into the algebraic dual of \(\mathcal{M}(\mathrm{T})\) , which extends the identity mapping of \(\mathcal{H}(\mathrm{T})\) . The relation \(\int f dm \in \mathrm{F} = \mathcal{H}(\mathrm{T})\) is therefore equivalent to \(f \in \mathcal{H}(\mathrm{T})\) .

Let \(u\) be a continuous linear mapping of \(F\) into a Hausdorff locally convex space \(G\) , and let us denote again by \(u\) its extension by bitransposition to a linear mapping of \(F'^*\) into \(G'^*\) (§1, No.~1). With this convention:

PROPOSITION 2. --- Every numerical function \(f\) essentially integrable for \(\mathbf{m}\) is essentially integrable for \(u \circ \mathbf{m}\) , and \(\int f d(u \circ \mathbf{m}) = u\left(\int f dm\right)\) .

The proposition is obvious, in view of the equality \(y'\circ u\circ m = {}^{t}u(y')\circ m\) for all \(y \in G'\) .

In general, if \(f \in \mathcal{L}(\mathbf{m})\) , the integral \(\int f dm\) belongs to \(\mathbf{F}'*\) but not to \(\mathbf{F}\) (see the above Example). However:

PROPOSITION 3. --- If the image under m of the set of \(f \in \mathcal{K}(\mathbf{T})\) such that \(\sup_{t \in \mathbf{T}} |f(t)| \leqslant 1\) is weakly relatively compact in F, then \(\int f dm \in \mathbf{F}\) for every bounded numerical function \(f\) essentially integrable for m.

Let \(\mathbf{A}\) be the set of \(f\in \mathcal{L}(\mathbf{m})\) such that \(\sup_{t\in \mathrm{T}}|f(t)|\leqslant 1\) , and let

\(B = A \cap \mathcal{K}(T)\) ; by hypothesis, \(\mathbf{m}(B)\) is weakly relatively compact in \(F\) , therefore it suffices to show that \(\mathbf{m}(A)\) is contained in the closure (in \(F'^*\) ) of \(\mathbf{m}(B)\) for the topology \(\sigma(F'^*, F')\) ; since \(\mathbf{m}(B)\) is convex and balanced, it suffices to prove that the polar of \(\mathbf{m}(\mathrm{B})\) in \(\mathrm{F}'\) is contained in that of \(\mathbf{m}(\mathrm{A})\) (TVS, II, §6, No.~3, Th. 1). Now, for a linear form \(\mathbf{z}' \in \mathrm{F}'\) to belong to \((\mathbf{m}(\mathrm{B}))^\circ\) , it is necessary and sufficient that \(|\langle \mathbf{z}', \mathbf{m}(g) \rangle| = |\int g d(\mathbf{z}' \circ \mathbf{m})| \leqslant 1\) for every function \(g \in \mathrm{B}\) , which signifies that the scalar measure \(|\mathbf{z}' \circ \mathbf{m}|\) is bounded and of norm \(\leqslant 1\) (Ch. III, §1, No.~8); but by (2) the latter condition implies that \(|\langle \mathbf{z}', \mathbf{m}(f) \rangle| \leqslant 1\) for every function \(f \in \mathrm{A}\) , whence \(\mathbf{z}' \in (\mathbf{m}(\mathrm{A}))^\circ\) .

COROLLARY 1. --- If, for every compact subset K of T, the image under m of the set of \(f \in \mathcal{K}(\mathrm{T},\mathrm{K})\) such that \(\sup_{t \in \mathrm{T}} |f(t)| \leqslant 1\) is weakly relatively compact in F, then \(\int f dm \in \mathrm{F}\) for every function \(f \in \mathcal{L}(\mathbf{m})\) that is bounded and has compact support, and \(\int f dm \in \mathrm{F}''\) for every function \(f \in \mathcal{L}(\mathbf{m})\) .

The first assertion may be deduced immediately from Prop. 3: if \(f\) is bounded and has compact support, and if \(U\) is a relatively compact open neighborhood of the support of \(f\) , then the restriction of \(\mathbf{m}\) to the subspace \(\mathcal{K}(U)\) is a measure \(\mathbf{m}_U\) on \(U\) that satisfies the conditions of Prop. 3, and \(\int f dm_U = \int f dm\) (Ch. V, §7, No.~1, Th. 1), therefore \(\int f dm \in F\) .

Now let \(f\) be any element of \(\mathcal{L}(\mathbf{m})\) ; for every compact subset \(K\) of \(T\) and every integer \(n > 0\) , let \(f_{n,K}\) be the numerical function on \(T\) defined as follows: if \(t \notin K\) , \(f_{n,K}(t) = 0\) ; if \(t \in K\) and \(|f(t) \leqslant n\) , \(f_{n,K}(t) = f(t)\) ; finally, if \(t \in K\) and \(|f(t)| > n\) , \(f_{n,K}(t) = nf(t)/|f(t)|\) . It is clear that for every \(t \in T\) , \(f(t)\) is the limit of \(f_{n,K}(t)\) with respect to the product filter of the Fréchet filter by the section filter of the (increasing directed) ordered set of compact subsets of \(T\) . Since \(|f_{n,K}| \leqslant |f|\) , it follows from Lebesgue's theorem and Prop. 10 of Ch. V, §1, No.~3, applied to each scalar measure \(|\mathbf{z}' \circ \mathbf{m}|\) , that \(f_{n,K}\) converges to \(f\) in \(\mathcal{L}(\mathbf{m})\) with respect to the preceding filter. Consequently, the integral \(\int f dm\) is in the closure in \(F'^*\) (for the topology \(\sigma(F'^*, F')\) ) of the set \(M\) of \(m(f_{n,K})\) . But the first part of the corollary shows that \(M \subset F\) , and, on the other hand, for every \(z' \in F'\) one has \(|\langle z', m(f_{n,K})\rangle| \leqslant \int |f| d|\mathbf{z}' \circ \mathbf{m}|\) , which shows that \(M\) is bounded in \(F_\sigma\) , hence also in \(F\) (TVS, IV, §1, No.~1, Prop. 1). Lemma 1 of §1, No.~2 therefore shows that \(\int f dm \in F''\) .

COROLLARY 2. --- If F is semi-reflexive, then \(\int f dm \in F\) for every numerical function \(f\) essentially integrable for \(\mathbf{m}\) .

